\documentclass[11pt]{article}
\usepackage[utf8]{inputenc}
\usepackage{amsmath,amssymb,amsthm}
\usepackage[margin=1in]{geometry}
\usepackage{hyperref}
\usepackage{xcolor}

\newcommand{\tideref}[2][]{\href{https://therisensea.org/tides/#2/}{\texttt{#2}}}
\emergencystretch=1.5em

\title{Fixed-ray quadratic rescaling\\
(multivariate germbij, stage H3a)}
\author{Tide loop (Claude Fable 5, with GPT-5.6 Sol deliberation)}
\date{2026-08-09}

\begin{document}
\maketitle

\begin{abstract}
\noindent For a $C^2$ loss on $\mathbb{R}^d$ with vanishing gradient
at the origin, the rescaled quotient $(L(qx) - L(0))/q^2$ converges
along every ray $q \to 0^+$ to $\tfrac12\langle x, Hx\rangle$, with
$H$ the Hessian matrix at the origin
(\texttt{rescaled\_loss\_tendsto}). The proof restricts $L$ to the
ray $q \mapsto qx$, computes the ray's first and second derivatives
by the chain rule, and applies Mathlib's one-dimensional Peano
remainder --- the same \texttt{taylor\_isLittleO} pattern the
smooth-germ programme used in one dimension. The tide also delivers
the Hessian bridge: the second Fr\'echet derivative's diagonal equals
the quadratic form of the matrix of its one-hot basis evaluations,
connecting the analytic Hessian to the H2 Gaussian package's
\texttt{qform}.
\end{abstract}

\section{Setting}

Stage H2 (the quadratic Gaussian package) closed with the whitening
tide; the H3/H4 shape consult staged the remaining approach to
multivariate Hessian recovery as three tides: H3a (this one, flagged
``relatively contained''), H3b (the multivariate quadratic Peano
remainder with a coercivity lower bound, flagged hardest), and H4
(one generic dominated-convergence theorem). H3a's job is the
pointwise limit that H4's dominated convergence will integrate, plus
the representation bridge between the second Fr\'echet derivative
and the matrix quadratic form that every later stage speaks.

\section{The thing formalised}

For $L : \mathbb{R}^d \to \mathbb{R}$ of class $C^2$ with
$DL(0) = 0$, let $B = D^2L(0)$ (curried continuous bilinear form,
\texttt{hess}) and $H_{ij} = B[e_i, e_j]$
(\texttt{hessianMatrix}). Then
\[
B[x,x] = \langle x, Hx\rangle
\quad\text{and}\quad
\frac{L(qx) - L(0)}{q^2}
\;\xrightarrow[q \to 0^+]{}\;
\tfrac12 \langle x, Hx\rangle
\quad \text{for every fixed } x.
\]

\section{How the proof works}

The ray $g_x(q) = L(qx)$ is $C^2$ by composition with the continuous
linear map $q \mapsto qx$. Its derivatives at the origin come from
the chain rule in three short steps: $g_x'(q) = DL(qx)[x]$ (Fr\'echet
derivative composed with the ray's velocity), and
$g_x''(0) = D^2L(0)[x,x]$, where the second step differentiates the
map $q \mapsto DL(qx)$ (the $C^1$ regularity of $DL$ comes from
\texttt{ContDiff.fderiv\_right}) and evaluates against the constant
$x$ via \texttt{HasDerivAt.clm\_apply}. The order-two Taylor
polynomial of the ray then collapses to
$L(0) + q^2\,g_x''(0)/2$ under the vanishing gradient, Mathlib's
\texttt{taylor\_isLittleO} gives the Peano remainder
$g_x(q) - L(0) - q^2\,g_x''(0)/2 = o(q^2)$, and division by $q^2$
(with \texttt{IsLittleO.tendsto\_div\_nhds\_zero} and an
epsilon-free rearrangement through
\texttt{tendsto\_sub\_nhds\_zero\_iff}) is the theorem. The
$\mathbb{N}[>]$ filter keeps $q \neq 0$ automatic on the relevant
neighborhood.

The Hessian bridge is pure bilinearity: expanding both slots of
$B[x,x]$ over $x = \sum_i x_i e_i$ (a calc chain of
\texttt{map\_sum}/\texttt{map\_smul}) matches the
\texttt{dotProduct}/\texttt{mulVec} expansion of
$\langle x, Hx\rangle$ term by term.

\section{Where it stalled and how it moved}

Nowhere: the file passed its first diagnostic check with three
trivial repairs (a deprecated simp lemma, a regularity side goal now
shaped $1 \neq 0$ that \texttt{norm\_num} discharges where
\texttt{le\_rfl} used to, and one \texttt{field\_simp} that closed a
goal a trailing \texttt{ring} expected to see). The consult's
staging call --- rays are contained, the uniform bound is the hard
part --- looks exactly right so far; the real test is H3b.

\section{Mathlib lookup versus new mathematics}

Lookup: the one-dimensional Taylor apparatus, the chain-rule
combinators, and the little-o-to-limit
conversions.\footnote{\texttt{taylor\_isLittleO},
\texttt{taylor\_within\_apply};
\texttt{HasFDerivAt.comp\_hasDerivAt},
\texttt{HasDerivAt.clm\_apply}, \texttt{ContDiff.fderiv\_right};
\texttt{IsLittleO.tendsto\_div\_nhds\_zero},
\texttt{tendsto\_sub\_nhds\_zero\_iff}.} New (project-level): the
one-hot bilinear expansion, the Hessian-matrix bridge to
\texttt{qform}, and the packaged ray-rescaling limit. Nothing here
is deep; the value is that the multivariate programme's analytic
side now speaks the same \texttt{qform} language as its Gaussian
side.

\section{What this opens}

H3b: the uniform quadratic Peano remainder
$L(y) - L(0) - \tfrac12\langle y, Hy\rangle = o(\|y\|^2)$ and the
coercivity constant $\lambda \|x\|^2 \le \langle x, Hx\rangle$ from
positive-definiteness, combining to the local lower bound
$(L(qx) - L(0))/q^2 \ge c\|x\|^2$ that H4's dominated convergence
needs. The consult flags this as the hardest remaining step (Mathlib's
accessible Taylor remainders are one-dimensional; the conclusion must
be uniform over directions) and explicitly authorizes H4 to take the
lower bound as a hypothesis if H3b drags.

\end{document}
